\section{Controlled and Observed Studies: Full Description}
\label{app:pollution_experiments}
\label{sec:evaluation}

We evaluate on a single New Delhi platform so the controlled and observed studies
share one regulatory sensor geometry, one PM\(_{2.5}\) and wind record, and one set
of proxy inventories.  We ask whether the geometry of \(\widetilde H_\Phi\) predicts
which coefficient, activity, and grouped-source claims hold.  A shared
base configuration is varied along one controlled axis at a time, avoiding
conflation of wind, source geometry, background flexibility, noise, and
forward-model error.  All runs use the \(40\times40\)
resolution, with a 72-hour transport window for the controlled experiments and
four consecutive one-week windows for the observed study.  Real imputed wind is a
controlled transport input when coefficients are synthetically known; the
observed-PM\(_{2.5}\) mode is a separate evaluation and is never assigned synthetic
source-recovery metrics.

\noindent\textbf{New Delhi Experimental Platform.}
The platform is built from hourly government PM\(_{2.5}\) and wind records
\citep{cpcb,cpcb_portal} spanning 1 May 2018--31 October 2020 over a 32-sensor
regulatory network.  PM\(_{2.5}\) is never imputed, and its valid-value mask defines
\(M_{\mathcal O}\) in Equation~\ref{eq:observation_selection_app}, with the same
ordered sensor--time rows retained in \(Y\), \(H_\Phi^{\mathrm{lag}}\), \(Q\), and
metadata.  Observed wind is converted (Equation~\ref{eq:wind_direction_conversion})
and completed on the response grid by the kernel coordinate-query imputer of
Appendix~\ref{subsec:wind_field_estimation}\label{subsec:new_delhi_wind}.  A learned
imputer was trained but not adopted (Appendix~\ref{subsec:results_wind_imputation}).
Four source groups (brick kilns, industries, population density, traffic) each
carry one normalized spatial proxy map and its own admissible temporal-basis
components (Table~\ref{tab:new_delhi_inventories},
Appendix~\ref{app:platform_details}), so fitted coefficients are in normalized
proxy units, not physical emission shares.  Both forward paths use the
open-boundary Gaussian puff operator of Appendix~\ref{subsec:plume_response}
evaluated at sensors, the matched generator for controlled recovery and the
inference operator for observed PM\(_{2.5}\).  An auxiliary advection--diffusion
simulator supplies a structural forward-model mismatch for the adequacy study
(Section~\ref{sec:identifiability_theory}), and normal backgrounds follow
Appendix~\ref{subsec:background_construction}.

The \emph{controlled} mode preserves this sensor geometry and transport while
assigning known nonnegative source--basis coefficients (synthetic or real gridded
wind, controlled observation noise, source splits, sensor-layout variants, and
optional operator mismatch), so it supports recovery metrics.  The \emph{observed} mode combines real PM\(_{2.5}\), its observation
mask, the imputed wind field, and the named normalized inventories; it has no
ground-truth source activities and is run over four consecutive one-week windows,
each fitted and reported independently.  Column-level records, Pusa-monitor
averaging and coverage, proxy normalization and crop, temporal-activity
assumptions, forward-model baselines, and simulator settings are given in
Appendix~\ref{app:platform_details}.\\


\noindent\textbf{Metrics and Reporting.}
Controlled trials, having known coefficients, report coefficient and
reconstructed-activity error, grouped activity/contribution error and projected
residual; interval coverage is not reported.  All trials also report the singular
spectrum and its diagnostics (\(\sigma_J\), numerical and effective rank, condition
number, visibility, background absorption, eligible coherence and ray distance), the
weak set, triggering pairs, and deterministic report groups
(Section~\ref{sec:identifiability_theory}).  They also record lag convergence, the
fixed-zero mask, background specification, and ensemble provenance (transport and
inventory kept separate).  Observed New Delhi runs instead report raw and
projected residuals, fitted trajectories, normalized proxy coefficients, group
signals in \(\mu\mathrm g/\mathrm m^3\), uncertainty intervals, inventory-scenario robustness,
weak flags, separations, group signals and shares, and residual calibration
status.  A reported share is an \(L_1\) fraction of the fitted projected group
signals, not a physical-emission share.

\subsection{Controlled Results}



\noindent\textbf{Recorded groupings.}  The controlled runs recorded groupings from the
pairwise rule (merge two declared blocks with a cross-source coherence above \(0.99\)).
On every saved operator, and in every observed week, this equals the reportable
partition at \(\tau_\theta=0.141\) with the most blocks (Section~\ref{sec:method}); for
Experiments 3 and 9 it also equals \(\mathcal P^\star\vee\mathcal D\).
\label{subsec:controlled_results}

The controlled trials have known source--basis coefficients and therefore recovery
metrics.  The controlled geometries are small (mostly
\(K=2\) source groups) so that conditioning, coherence, and background absorption
can each be varied in isolation and read directly off \(\widetilde H_\Phi\).  The
observed study then exercises the full four-group, seven-coefficient pipeline on
real data.  Figure~\ref{fig:controlled} collects the baseline comparison and the
seven controlled diagnostics.  Full numeric tables are in
Appendix~\ref{app:expanded_results}.


\noindent\textbf{Baselines.}
\label{subsec:results_baselines}
We compare IASA with plain NNLS on the projected system (the same estimator without
grouping or diagnostics), chemical mass balance on the unprojected response, and a
PMF/NMF receptor factorization, over three seeds (Figure~\ref{fig:controlled}(a);
Table~\ref{tab:results_baselines}).  IASA raises an identifiability flag in every
scenario and seed and none of the baselines does.  Under background stress the
projected source signals stay separated; IASA's largest group-signal error is
\(0.028\)--\(0.060\) against \(2.10\)--\(4.63\) for chemical mass balance, and plain NNLS
matches IASA.  In the steady-wind case, two seeds have no signal at the sensors; in the third
IASA merges the two sources and the merged group's signal error is \(0.004\), against
\(7.84\) for chemical mass balance, while the baselines report the individual split
without a flag.

\noindent\textbf{Conditioning and Recovery.}
\label{subsec:results_h1}
We vary source geometry and Gaussian observation noise (0--20\% of the maximum
clean sensor signal).  Both geometries are full rank (numerical rank $=$ effective
rank $=2$), so the fixed geometry fixes $\sigma_J$ and $\kappa$ and coefficient
error grows linearly with noise.  The better-conditioned \emph{close} geometry
recovers proportionally more accurately than the \emph{separated} one
(Figure~\ref{fig:controlled}(b); Table~\ref{tab:results_h1}): at equal noise the
coefficient error is set by the response-matrix geometry.

\noindent\textbf{Wind Diversity and Sensor Geometry.}
\label{subsec:results_h4}
Matched source, basis, and sensor sets are evaluated across six wind providers
(constant, single-direction, diurnal, AR(1), multi-directional, and real gridded
New Delhi wind) and three sensor layouts (regulatory, seeded random,
downwind-focused), with identical source--basis columns on both sides of every
comparison (Figure~\ref{fig:controlled}(c); Table~\ref{tab:results_h4}).  Two
effects are seen: wind diversity improves conditioning (multi-directional and real
wind give the lowest maximum eligible coherence, while a single steady direction
pushes it toward one), and the sensor layout changes $\sigma_J$ as much as the wind
does: with the seeded random layout one source is nearly invisible, $\sigma_J$ is near
zero and the only nonzero coefficient errors in the sweep occur, while the regulatory
layout keeps $\sigma_J\ge1.53$.  Separate historical and AR(1) window ensembles
both attain full rank with probability $1.0$ and no pairwise ambiguity
(Table~\ref{tab:results_h4_ens}).

\noindent\textbf{Background Stress.}
\label{subsec:results_h3}
We compare no background, the normal rank-four basis, a redundant-column basis with
the same span, and a labeled source-like stress basis, all declared before fitting.
The redundant basis is byte-identical to the primary in every reported quantity,
while the source-like \emph{stress} basis drives $\sigma_J$ and minimum visibility
to $0$ and absorption to $1$ (Figure~\ref{fig:controlled}(d);
Table~\ref{tab:results_h3}).  A background basis that can reproduce a source's
signal absorbs it and removes its identifiability while the fitted residual is
unchanged.  The coefficient error \emph{falls} in the stress case because the signal
has been absorbed.

\noindent\textbf{Transport Error.}
Parametric perturbations of wind speed, direction, and dispersion within the puff
family propagate into attribution.  Coefficient error grows with the operator-error
norm along the direction (to $0.86$ at $20^\circ$) and dispersion (to $0.89$ at
$2\times$) axes, while the speed axis is non-monotonic, so the operator-error norm does not order
the attribution error (Figure~\ref{fig:controlled}(e);
Table~\ref{tab:results_h5a}).  A structural mismatch (open-boundary puff fitted to
edge-hold advection--diffusion observations, mismatch norm $0.44$) drives mean
coefficient error to $0.86$ and the refitted-bootstrap adequacy test rejects in
$100\%$ of trials, so the test also rejects forward-model misspecification, not only
missing sources.

\noindent\textbf{Lag-Window Selection.}
Physical travel and residence times define an increasing candidate lag grid.  As the
window grows, conditioning improves monotonically
($\kappa$ falls from $569$ at $L{=}4$ to $5.1$ at $L{=}16$; $\sigma_J$ rises from
$0.08$ to $9.07$) while numerical rank and report-component count stay fixed and
recovery is exact (Figure~\ref{fig:controlled}(f); Table~\ref{tab:results_lag}).
The relative Frobenius changes of Eq.~\ref{eq:lag_convergence} between adjacent
candidates are $0.063$, $0.096$, $0.082$, $0.057$ and $0.139$, so no candidate meets
$\tau_L=10^{-3}$ and the convergence rule does not select a lag on this grid; the
other reported runs use the declared $L=12$.

\noindent\textbf{Temporal-Basis Recovery.}
Reconstructing known diurnal, intermittent, and day/night source--basis
coefficients at increasing noise separates coefficient error from
reconstructed-activity error.  The activity-trajectory error stays consistently below
the coefficient error ($0.58$ vs.\ $0.90$ at $20\%$ noise), at a ratio of $0.57$ up to
$10\%$ noise and $0.65$ at $20\%$ (Figure~\ref{fig:controlled}(g); Table~\ref{tab:results_temporal}),
because projecting noisy coefficients through the temporal bases damps isotropic
estimation noise more than the coherent signal.

\noindent\textbf{Inventory Robustness.}
Perturbing inventory locations, spatial scales, category assignments, and map
versions (each a separate scenario, never pooled with transport error) recovers the
known coefficients exactly, but the identifiability geometry moves with the
inventory version.  Rescaling or substituting the map multiplies $\sigma_J$ by $2.18$ and $2.13$ (to
$47.0$ and $46.0$ vs.\ $21.6$ baseline) while a symmetric category swap leaves it
unchanged (Figure~\ref{fig:controlled}(h); Table~\ref{tab:results_h5b}).  The
attribution \emph{target} thus changes with the assumed inventory even when recovery
is exact.  The tables for coherence and grouped reporting, missing-source adequacy
and per-sensor footprints are in Appendix~\ref{app:additional_experiments}.


\section{Expanded Results}
\label{app:expanded_results}

Full numeric tables backing Figures~\ref{fig:controlled} and~\ref{fig:observed}.

\noindent\textbf{Summary of further controlled results.}
From \(L=4\) to \(16\) hours \(\sigma_J\) rises from \(0.08\) to \(9.07\), but the operator
still changes by \(6\)--\(14\%\) between adjacent lags (Experiment 7); activity
trajectories have lower error than coefficients (Experiment 9); \(\sigma_J\) moves from
\(21.6\) to \(47.0\) under a rescaled inventory map although recovery stays exact
(Experiment 6); and a shifted copy of one map raises the coherence to \(0.958\)
(separation \(0.286\)) without a merge (Experiment 2)
(Figure~\ref{fig:controlled}(f)--(h), Appendix~\ref{app:expanded_results}).
Eight downwind sensors give coherence \(0.9897\) and \(0.9901\) under constant and
single-direction wind (Table~\ref{tab:results_h4}), on either side of \(0.99\), so the
rule keeps the two sources apart in the first case and merges them in the second (the
run records coherences, not groupings).

\subsection{Conditioning (Exp 1)}
Coefficient recovery versus conditioning and noise; the better-conditioned geometry
recovers more accurately at matched noise.
\begin{table}[t]
\centering
\footnotesize
\caption{Experiment 1: coefficient recovery versus conditioning and noise. The
better-conditioned geometry recovers more accurately at matched noise.}
\label{tab:results_h1}
\begin{tabular}{llrrrr}
\toprule
geom. & noise & $\sigma_J$ & $\kappa$ & coef.\ err. & resid. \\
\midrule
separated & 0.00 & 3.27 & 8.89 & 0.000 & 0.00 \\
separated & 0.05 & 3.27 & 8.89 & 0.066 & 6.37 \\
separated & 0.10 & 3.27 & 8.89 & 0.132 & 12.73 \\
separated & 0.20 & 3.27 & 8.89 & 0.264 & 25.46 \\
close     & 0.00 & 9.60 & 2.26 & 0.000 & 0.00 \\
close     & 0.05 & 9.60 & 2.26 & 0.028 & 8.42 \\
close     & 0.10 & 9.60 & 2.26 & 0.057 & 16.84 \\
close     & 0.20 & 9.60 & 2.26 & 0.113 & 33.68 \\
\bottomrule
\end{tabular}
\end{table}

\subsection{Wind Diversity and Sensor Geometry (Exp 4)}
$\sigma_J$, maximum eligible coherence, and coefficient error across wind provider
$\times$ sensor layout, and the wind-window ensemble spread.  With the eight randomly
placed sensors and the constant, single, diurnal and AR(1) winds, the operator column
of source B has norm $5.2\times10^{-7}$, $1.7\times10^{-4}$, $5.2\times10^{-5}$ and
$1.1\times10^{-4}$, against $7.4\times10^{-4}$, $0.44$, $0.065$ and $0.22$ for
source A and $7.6$--$18.6$ for both sources under the regulatory layout (recomputed
from the same configuration).  In these cells $\sigma_J$ is near zero because source B
is nearly invisible, and the coherence is computed on a column that carries almost no
signal.
\begin{table}[t]
\centering
\footnotesize
\caption{Experiment 4: $\sigma_J$, maximum eligible coherence, and coefficient
error across wind provider $\times$ sensor layout.}
\label{tab:results_h4}
\begin{tabular}{llrrr}
\toprule
wind & layout & $\sigma_J$ & max.\ coh. & coef.\ err. \\
\midrule
constant & regulatory & 1.53 & 0.975 & 0.000 \\
constant & random     & 0.00 & 0.028 & 0.000 \\
constant & downwind   & 0.13 & 0.9897 & 0.000 \\
single   & regulatory & 6.44 & 0.741 & 0.000 \\
single   & random     & 0.00 & 1.000 & 0.573 \\
single   & downwind   & 1.36 & 0.9901 & 0.000 \\
diurnal  & regulatory & 6.11 & 0.722 & 0.000 \\
diurnal  & random     & 0.00 & 0.880 & 0.001 \\
diurnal  & downwind   & 2.60 & 0.897 & 0.000 \\
ar1      & regulatory & 7.35 & 0.641 & 0.000 \\
ar1      & random     & 0.00 & 0.774 & 0.115 \\
ar1      & downwind   & 2.14 & 0.916 & 0.000 \\
multi    & regulatory & 5.70 & 0.129 & 0.000 \\
multi    & random     & 2.79 & 0.322 & 0.000 \\
multi    & downwind   & 0.41 & 0.110 & 0.000 \\
real     & regulatory & 2.85 & 0.156 & 0.000 \\
real     & random     & 2.10 & 0.287 & 0.000 \\
real     & downwind   & 0.38 & 0.221 & 0.000 \\
\bottomrule
\end{tabular}
\end{table}
\begin{table}[t]
\centering
\footnotesize
\caption{Experiment 4: wind-window ensemble distribution ($\sigma_J$ percentiles;
both ensembles reach full rank with probability 1 and show no ambiguity).}
\label{tab:results_h4_ens}
\begin{tabular}{lrrr}
\toprule
ensemble & $\sigma_J^{5}$ & $\sigma_J^{50}$ & $\sigma_J^{95}$ \\
\midrule
historical (real slices) & 1.45 & 4.51 & 9.96 \\
simulated (AR(1))        & 3.50 & 5.58 & 6.80 \\
\bottomrule
\end{tabular}
\end{table}

\subsection{Background Stress (Exp 3)}
Redundant $\equiv$ primary (same span); the source-like stress basis drives
$\sigma_J$ to $0$ and absorption to $1$.
\begin{table}[t]
\centering
\footnotesize
\caption{Experiment 3: background stress. Redundant $\equiv$ primary (same span);
the source-like stress basis drives $\sigma_J$ to $0$ and absorption to $1$.}
\label{tab:results_h3}
\begin{tabular}{lrrrr}
\toprule
background & min.\ vis. & absorp. & $\sigma_J$ & coef.\ err. \\
\midrule
none      & 1.000 & 0.000 & 0.080 & 2.36 \\
primary   & 0.972 & 0.235 & 0.080 & 2.61 \\
redundant & 0.972 & 0.235 & 0.080 & 2.61 \\
stress    & 0.000 & 1.000 & 0.000 & 0.78 \\
\bottomrule
\end{tabular}
\end{table}

\subsection{Transport Error (Exp 5)}
\label{subsec:results_h5a}
Parametric perturbations (wind speed, direction, dispersion); a structural PDE
mismatch generated with the auxiliary edge-hold advection--diffusion solver yields
adequacy rejection rate $1.0$.
\begin{table}[t]
\centering
\footnotesize
\caption{Experiment 5: parametric transport error. Coefficient error rises with
the operator error norm on the direction and dispersion axes (the speed axis is
non-monotonic). Structural PDE mismatch: adequacy rejection rate $1.0$.}
\label{tab:results_h5a}
\begin{tabular}{llrrr}
\toprule
axis & value & op.\ err. & $\sigma_J$ & coef.\ err. \\
\midrule
direction ($^\circ$) & 0  & 0.00 & 0.87 & 0.00 \\
direction ($^\circ$) & 5  & 0.29 & 3.37 & 0.62 \\
direction ($^\circ$) & 10 & 0.60 & 7.07 & 0.80 \\
direction ($^\circ$) & 20 & 0.87 & 7.98 & 0.86 \\
speed ($\times$)     & 1.25 & 0.29 & 2.09 & 0.45 \\
speed ($\times$)     & 1.50 & 0.53 & 4.57 & 0.27 \\
dispersion ($\times$)& 1.50 & 0.19 & 1.42 & 0.50 \\
dispersion ($\times$)& 2.00 & 0.35 & 1.76 & 0.89 \\
\bottomrule
\end{tabular}
\end{table}

\subsection{Inventory Robustness (Exp 6)}
\label{subsec:results_h5b}
Location, scale, category, and map-version perturbations, each a separate scenario
never pooled with transport error; recovery exact, $\sigma_J$ tracks the version.
\begin{table}[t]
\centering
\footnotesize
\caption{Experiment 6: inventory robustness on the real maps. Recovery is exact in
every scenario; $\sigma_J$ tracks the inventory version.}
\label{tab:results_h5b}
\begin{tabular}{lrr}
\toprule
scenario & $\sigma_J$ & coef.\ err. \\
\midrule
baseline        & 21.56 & 0.0 \\
location shift   & 28.85 & 0.0 \\
spatial rescale  & 47.00 & 0.0 \\
alt.\ map version & 46.01 & 0.0 \\
category swap     & 21.56 & 0.0 \\
\bottomrule
\end{tabular}
\end{table}

\subsection{Lag-Window Selection (Exp 7)}
\label{subsec:results_lag_selection}
Conditioning improves with lag.  No candidate meets the convergence threshold
$\tau_L=10^{-3}$ of Equation~\ref{eq:lag_convergence}; the code then records the largest
candidate ($L=16$) as a fallback.  The other reported runs use the declared $L=12$.
\begin{table}[t]
\centering
\footnotesize
\caption{Experiment 7: lag window.  Conditioning improves with lag; no candidate meets
$\tau_L$, and the other reported runs use the declared $L=12$.}
\label{tab:results_lag}
\begin{tabular}{rrrrr}
\toprule
lag $L$ & $\sigma_J$ & $\kappa$ & \#comp. & coef.\ err. \\
\midrule
4  & 0.08 & 569.2 & 2 & 0.0 \\
6  & 0.71 & 63.8  & 2 & 0.0 \\
8  & 2.46 & 18.6  & 2 & 0.0 \\
10 & 4.57 & 10.1  & 2 & 0.0 \\
12 & 6.16 & 7.5   & 2 & 0.0 \\
16 & 9.07 & 5.1   & 2 & 0.0 \\
\bottomrule
\end{tabular}
\end{table}

\subsection{Temporal-Basis Recovery (Exp 9)}
\label{subsec:results_temporal_basis}
Activity-trajectory error stays below coefficient error at every noise level (fixed
$\sigma_J$), a fixed geometric property of the temporal-basis map.
\begin{table}[t]
\centering
\footnotesize
\caption{Experiment 9: temporal-basis recovery (bases: diurnal, block, day/night).
Activity-trajectory error stays below coefficient error at every noise level.}
\label{tab:results_temporal}
\begin{tabular}{rrr}
\toprule
noise & coef.\ err. & activity err. \\
\midrule
0.00 & 0.000 & 0.000 \\
0.02 & 0.132 & 0.075 \\
0.05 & 0.329 & 0.188 \\
0.10 & 0.659 & 0.375 \\
0.20 & 0.900 & 0.582 \\
\bottomrule
\end{tabular}
\end{table}

\subsection{Baseline Comparison (Exp 11)}
IASA against plain NNLS on the projected system (the same estimator without grouping
or diagnostics), chemical mass balance (NNLS on the unprojected response) and a
PMF/NMF receptor factorization, on the steady-wind case (one wind direction, six random sensors) and the source-like
background, seeds \(0\)--\(2\), run on CPU.  Shares are \(L_1\) fractions of the
projected group signals \(\At_k\widehat{\mathbf c}_k\), the quantity the signal target
identifies; the group-signal error is
\(\|\At_G(\widehat{\mathbf c}_G-\mathbf c_G)\|_2/\|\At_G\mathbf c_G\|_2\) for the blocks
of IASA's reported partition, and the table gives the largest over blocks.  PMF/NMF has
no coefficients in the operator's units, so only its share error on the unprojected
response is given.  In seeds \(0\) and \(1\) of the steady-wind case both operator columns have
norm below \(10^{-10}\) (against \(0.06\) in seed \(2\)): no signal reaches the six
sensors, IASA flags \(\sigma_J\approx0\), and no error is reported.  Shares computed
from the unprojected response include the absorbed column of the background-stress
scenario, whose coefficient is not identified: plain NNLS assigns it values above
\(10^{13}\) in seeds \(1\) and \(2\), which gives unprojected share errors of \(1.30\) and
\(1.26\) although its group-signal errors are within \(0.003\) of IASA's.
\begin{table}[t]
\centering
\footnotesize
\setlength{\tabcolsep}{3pt}
\caption{Experiment 11, seeds \(0,1,2\): projected share error (\(L_2\)) and largest
relative group-signal error.  In the steady-wind case only seed \(2\) has signal at the
sensors; IASA merges its two sources and reports the merged group.  \(^\dagger\)Share
error on the unprojected response.}
\label{tab:results_baselines_full}
\begin{tabular}{lllll}
\toprule
scenario & method & flag & share err. & group-signal err. \\
\midrule
bg.\ stress & IASA       & yes & \(0.007,0.003,0.007\) & \(0.054,0.028,0.060\) \\
bg.\ stress & plain NNLS & no  & \(0.007,0.003,0.006\) & \(0.054,0.028,0.062\) \\
bg.\ stress & CMB        & no  & \(0.193,0.123,0.190\) & \(4.63,2.10,4.32\) \\
bg.\ stress & PMF/NMF    & no  & \(1.15,1.19,1.18\)\(^\dagger\) & -- \\
\midrule
steady wind (s.\ 2) & IASA       & yes & merged & \(0.004\) \\
steady wind (s.\ 2) & plain NNLS & no  & \(0.036\) & \(0.004\) \\
steady wind (s.\ 2) & CMB        & no  & \(0.558\) & \(7.84\) \\
steady wind (s.\ 2) & PMF/NMF    & no  & \(0.304\)\(^\dagger\) & -- \\
\bottomrule
\end{tabular}
\end{table}

\subsection{Observed Weeks: Full Analysis}
\label{app:observed_full}

\noindent\textbf{Data and windows.}
We use hourly government PM\(_{2.5}\) and wind records \citep{cpcb,cpcb_portal} from
the 32-sensor New Delhi regulatory network, four proxy source groups (brick kilns,
industries, population density, traffic) with declared temporal bases, and the
gridded wind field of Appendix~\ref{app:application}.  After the fixed-zero set,
\(J=7\): one temporal component each for brick kilns, industries and population, and
four time-of-day components for traffic.  The declared partition \(\mathcal D\) has
one block per source group.  Four consecutive one-week windows, 1--28 May 2018, are
fitted and reported separately, each on the \(79\)--\(85\%\) of sensor-hours with a
valid reading (\(N=4{,}236\)--\(4{,}561\)).  There is no ground truth for source
contributions.

\noindent\textbf{Noise level.}
The two co-located Pusa monitors differ with standard deviation
\(43.4\,\mu\mathrm g/\mathrm m^3\) over \(17{,}949\) shared hours (none in May 2018), so
if their errors are independent with equal variance each has
\(\sigma_e=30.7\,\mu\mathrm g/\mathrm m^3\).  The difference has lag-1 autocorrelation
\(0.857\); for errors independent across sensors and AR(1) in time with this
coefficient, the error covariance has spectral norm at most
\((1+0.857)/(1-0.857)=13.0\) times \(\sigma_e^2\), an effective noise level of
\(110.7\,\mu\mathrm g/\mathrm m^3\).  Both describe instrument error only; operator
error adds to them (Proposition~\ref{prop:operator_perturbation}).



\noindent\textbf{Results.}
\label{subsec:results_new_delhi_identifiability}
\label{subsec:results_new_delhi_apportionment}
Section~\ref{subsec:observed_results} gives the identifiable resolution, the group
signals against the noise level and the shares with their intervals.  The tables
below give the per-week identifiability, the residual diagnostics and, for each
source group, its separation (the sine of the smallest angle between the span of the
group's fingerprints and the span of the other fingerprints) and the root mean square
of its fitted signal per reading.  The fit subtracts each week's transported
initial-condition baseline, and the residual adequacy test is reported as
uncalibrated because no calibrated noise model for the observed PM\(_{2.5}\) is
available.

\subsection{Observed New Delhi (weeks 1--4)}
Per-week identifiability geometry and sensor-fit/residual diagnostics (PM$_{2.5}$
never imputed; adequacy uncalibrated).
\begin{table}[t]
\centering
\footnotesize
\caption{Observed New Delhi identifiability, weeks 1--4. All weeks are full rank,
finite-conditioned, reported as the four source groups (traffic has four columns).}
\label{tab:results_nd_ident}
\begin{tabular}{lrrrrl}
\toprule
week & $\sigma_1$ & $\sigma_J$ & num.\ rank & max.\ coh. & groups \\
\midrule
1 & 50.1 & 3.71  & 7 & 0.316 & 4 groups \\
2 & 60.3 & 5.67  & 7 & 0.355 & 4 groups \\
3 & 77.0 & 10.15 & 7 & 0.496 & 4 groups \\
4 & 81.2 & 7.65  & 7 & 0.475 & 4 groups \\
\bottomrule
\end{tabular}
\end{table}
\begin{table}[t]
\centering
\footnotesize
\caption{Observed New Delhi sensor fit and residual diagnostics, weeks 1--4.
PM$_{2.5}$ never imputed; adequacy uncalibrated (no verdict).}
\label{tab:results_nd_resid}
\begin{tabular}{lrrrl}
\toprule
week & obs.\ rows & mask frac. & resid.\ norm & adequacy \\
\midrule
1 & 4561 & 0.848 & 3768 & uncalibrated \\
2 & 4278 & 0.796 & 5955 & uncalibrated \\
3 & 4315 & 0.803 & 3050 & uncalibrated \\
4 & 4236 & 0.788 & 3760 & uncalibrated \\
\bottomrule
\end{tabular}
\end{table}

\begin{table}[t]
\centering
\footnotesize
\setlength{\tabcolsep}{2.5pt}
\caption{Observed weeks, largest group: share of the fitted projected signal, 95\%
bootstrap interval and \(b_g/\|\At_g\widehat{\mathbf c}_g\|_2\) at the instrument noise
level \(30.7\) and the AR(1) upper level \(110.7\,\mu\mathrm g/\mathrm m^3\).  All groups:
Table~\ref{tab:observed_shares_full}.  The bootstrap draws independent noise at each
level; its percentile intervals need not contain the estimate, because NNLS sets some
coefficients to zero.}
\label{tab:observed_shares}
\vspace{3pt}
\begin{tabular}{llrcccc}
\toprule
 & & & \multicolumn{2}{c}{interval} & \multicolumn{2}{c}{bound/norm} \\
\cmidrule(lr){4-5}\cmidrule(lr){6-7}
wk & group & share & \(30.7\) & \(110.7\) & \(30.7\) & \(110.7\) \\
\midrule
1 & population & 1.00 & 0.75--1.00 & 0.26--0.99 & 0.49 & 1.78 \\
2 & population & 0.81 & 0.71--0.85 & 0.49--0.93 & 0.24 & 0.88 \\
3 & population & 0.94 & 0.00--0.94 & 0.00--0.92 & 2.96 & 10.67 \\
4 & brick kilns & 0.93 & 0.33--1.00 & 0.00--1.00 & 1.33 & 4.81 \\
\bottomrule
\end{tabular}
\end{table}

Table~\ref{tab:results_nd_signal} gives the signal-target analysis of Section~\ref{subsec:observed_results}, computed from the committed projected operators and fits of the four weeks with \texttt{experiments/iasa\_pol/signal\_target\_weeks.py}.  The bootstrap adds Gaussian noise in \(\operatorname{col}(\At)\) to the fitted signal, which is the only part of the noise that changes the NNLS fit, and refits with an active-set NNLS solver (\texttt{scipy.optimize.nnls}), whose solution agrees with the projected-FISTA fit to five decimals on the four weeks.  The reported groupings in this table and in the Experiment 11 scoring were computed with a separation rule (every block must have separation at least \(\sqrt{1-0.99^2}=0.141\)); on these operators it gives the same blocks as the pairwise coherence rule.

\begin{table}[t]
\centering
\footnotesize
\setlength{\tabcolsep}{2.5pt}
\caption{Observed New Delhi on the signal target.  \(s_g\): separation of the source
group; signal: root mean square of the fitted projected group signal per reading
(\(\mu\mathrm g/\mathrm m^3\)); share with 95\% parametric-bootstrap interval
at \(\sigma_e=30.7\) (\(1000\) refits).}
\label{tab:results_nd_signal}
\begin{tabular}{llrrrl}
\toprule
wk & group & \(s_g\) & signal & share & interval \\
\midrule
1 & brick kilns & 0.89 & 0.00 & 0.00 & 0.00--0.09 \\
 & industries & 0.95 & 0.00 & 0.00 & 0.00--0.14 \\
 & population & 0.87 & 5.40 & 1.00 & 0.75--1.00 \\
 & traffic & 0.82 & 0.00 & 0.00 & 0.00--0.17 \\
\midrule
2 & brick kilns & 0.94 & 0.00 & 0.00 & 0.00--0.04 \\
 & industries & 0.93 & 2.51 & 0.16 & 0.10--0.21 \\
 & population & 0.83 & 11.86 & 0.81 & 0.71--0.85 \\
 & traffic & 0.84 & 0.96 & 0.04 & 0.01--0.10 \\
\midrule
3 & brick kilns & 0.96 & 0.09 & 0.06 & 0.00--0.45 \\
 & industries & 0.86 & 0.00 & 0.00 & 0.00--0.58 \\
 & population & 0.78 & 1.03 & 0.94 & 0.00--0.94 \\
 & traffic & 0.87 & 0.00 & 0.00 & 0.00--0.88 \\
\midrule
4 & brick kilns & 0.96 & 1.87 & 0.93 & 0.33--1.00 \\
 & industries & 0.87 & 0.08 & 0.07 & 0.00--0.41 \\
 & population & 0.72 & 0.00 & 0.00 & 0.00--0.32 \\
 & traffic & 0.77 & 0.00 & 0.00 & 0.00--0.55 \\
\bottomrule
\end{tabular}
\end{table}

\begin{table}[t]
\centering
\footnotesize
\setlength{\tabcolsep}{4pt}
\caption{Observed weeks, all groups: share of the fitted projected signal and 95\%
parametric-bootstrap intervals (\(1{,}000\) refits) at the instrument noise level
\(\sigma_e=30.7\) and at the AR(1) upper level \(110.7\) (both in
\(\mu\mathrm g/\mathrm m^3\)); the bootstrap draws independent noise at each level.}
\label{tab:observed_shares_full}
\vspace{3pt}
\begin{tabular}{llrcc}
\toprule
wk & group & share & at \(30.7\) & at \(110.7\) \\
\midrule
1 & brick kilns & 0.00 & 0.00--0.09 & 0.00--0.26 \\
 & industries & 0.00 & 0.00--0.14 & 0.00--0.45 \\
 & population & 1.00 & 0.75--1.00 & 0.26--0.99 \\
 & traffic & 0.00 & 0.00--0.17 & 0.00--0.51 \\
\midrule
2 & brick kilns & 0.00 & 0.00--0.04 & 0.00--0.11 \\
 & industries & 0.16 & 0.10--0.21 & 0.00--0.33 \\
 & population & 0.81 & 0.71--0.85 & 0.49--0.93 \\
 & traffic & 0.04 & 0.01--0.10 & 0.00--0.27 \\
\midrule
3 & brick kilns & 0.06 & 0.00--0.45 & 0.00--0.67 \\
 & industries & 0.00 & 0.00--0.58 & 0.00--0.86 \\
 & population & 0.94 & 0.00--0.94 & 0.00--0.92 \\
 & traffic & 0.00 & 0.00--0.88 & 0.00--1.00 \\
\midrule
4 & brick kilns & 0.93 & 0.33--1.00 & 0.00--1.00 \\
 & industries & 0.07 & 0.00--0.41 & 0.00--0.74 \\
 & population & 0.00 & 0.00--0.32 & 0.00--0.61 \\
 & traffic & 0.00 & 0.00--0.55 & 0.00--1.00 \\
\bottomrule
\end{tabular}
\end{table}

\section{Additional Experiments}
\label{app:additional_experiments}

Three further controlled experiments not shown in the main-text figures.

\subsection{Coherence and Grouped Reporting (Exp 2)}
\label{subsec:results_h2}
We form increasingly coherent source pairs by shifting a copy of one inventory map by
a controlled spatial offset, so decreasing the shift drives coherence toward one.
Shrinking the offset drives maximum eligible coherence from $0.05$ to $0.96$ while
$\sigma_J$ falls from $21.8$ to $4.6$, yet at every offset the two sources stayed separate groups with exact recovery
(individual relative error $0$; Table~\ref{tab:results_h2}).  The largest coherence,
$0.958$, is below $\tau_\rho=0.99$, so the pairwise rule keeps the two sources apart.
\begin{table}[t]
\centering
\footnotesize
\caption{Experiment 2: coherence rises as the source offset shrinks; no pair exceeds
\(\tau_\rho=0.99\) and noiseless recovery stays exact.}
\label{tab:results_h2}
\begin{tabular}{rrrll}
\toprule
offset & max.\ coh. & $\sigma_J$ & merged & indiv./grp.\ err. \\
\midrule
8.0 & 0.048 & 21.77 & no & 0.0 / 0.0 \\
6.0 & 0.862 & 10.48 & no & 0.0 / 0.0 \\
4.0 & 0.905 & 8.40  & no & 0.0 / 0.0 \\
2.0 & 0.945 & 5.43  & no & 0.0 / 0.0 \\
1.0 & 0.958 & 4.64  & no & 0.0 / 0.0 \\
\bottomrule
\end{tabular}
\end{table}

\subsection{Missing-Source Adequacy (Exp 8)}
\label{subsec:results_missing_source}
On the real New Delhi platform with a non-empty rank-4 primary background, the
refitted parametric bootstrap ($\alpha=0.05$, $1000$ replicates) rejects at rate
$0.05$ under the null and at rate $1.0$ against a residual-visible omitted
source whose signature lies largely outside
$\operatorname{span}[H_\Phi^{\mathrm{lag}},Q]$ (rejection $1.0$); an \emph{aligned}
omission, a nonnegative combination of the fitted columns and background, is absorbed
and not detected (rejection $0.05$; Table~\ref{tab:results_missing}). A rejection is
therefore evidence of model inadequacy, whereas a non-rejection does not show that the
inventory is complete.
\begin{table}[t]
\centering
\footnotesize
\caption{Experiment 8: missing-source adequacy on the platform (non-empty $Q$).
Calibrated null, full power on a residual-visible omission, and no detection of an
absorbed in-span omission.}
\label{tab:results_missing}
\begin{tabular}{lr}
\toprule
omission case & rejection rate \\
\midrule
none (null, calibration) & 0.05 \\
residual-visible (out of span, frac.\ 0.91) & 1.00 \\
aligned (in span, absorbed) & 0.05 \\
\bottomrule
\end{tabular}
\end{table}

\subsection{Per-Sensor Footprints (Exp 10)}
\label{subsec:results_footprints}
On controlled trials the per-sensor footprints of
Appendix~\ref{app:footprint_details} localize the responsible upwind cells (mass
centroid $0.74$ cells from the true origin, $100\%$ of mass within the $4$-cell
radius, nonnegative) and per-sensor contributions sum to the fitted sensor signal to
machine precision; because a non-singleton component is present, footprints are
reported at group granularity. On observed New Delhi data footprints are reported per
monitor at the identifiable report groups, accompanying the apportionment of
Section~\ref{subsec:observed_results}
(Appendix~\ref{app:footprint_details}).

